\section{Khảo sát và đánh giá các giải pháp hỗ trợ quyết định lâm sàng hiện hữu}
\label{sec:existing_solutions}

Trên thế giới và tại Việt Nam, đã có một số hệ thống và phần mềm hỗ trợ kiểm tra tương tác thuốc được đưa vào ứng dụng:

\begin{table}[htbp]
\centering
\small
\begin{tabularx}{\textwidth}{L{3.0cm} L{3.5cm} L{3.5cm} Y}
\toprule
\textbf{Hệ thống} & \textbf{Ưu điểm nổi bật} & \textbf{Hạn chế kỹ thuật} & \textbf{Khả năng áp dụng tại Việt Nam} \\
\midrule
\textbf{UpToDate Lexicomp} & Cơ sở dữ liệu y văn thế giới khổng lồ, cập nhật liên tục các nghiên cứu lâm sàng quốc tế. & Chi phí bản quyền rất đắt đỏ (hàng nghìn USD/năm); hoàn toàn không có danh mục thanh toán BHYT Việt Nam; giao diện tiếng Anh. & Phù hợp cho bệnh viện tư nhân lớn hoặc nghiên cứu; không tối ưu cho vận hành đại trà bệnh viện công. \\
\addlinespace
\textbf{Cerner Multum / Epic CDSS} & Tích hợp sâu vào hệ thống Bệnh án Điện tử (EMR) hàng đầu thế giới; khả năng phân tích đa bệnh lý mạnh. & Hệ thống cồng kềnh, chi phí triển khai hàng triệu USD; đòi hỏi hạ tầng máy chủ trung tâm phức tạp; không có dữ liệu DĐVN. & Khó tiếp cận đối với đại đa số bệnh viện công lập tại Việt Nam. \\
\addlinespace
\textbf{Mô-đun CDSS trong HIS nội địa (VNPT, Viettel, FPT)} & Tích hợp sẵn trong phần mềm quản lý bệnh viện; hỗ trợ danh mục BHYT và xuất hóa đơn viện phí. & Thuật toán kiểm tra luật tĩnh đơn giản; tỷ lệ cảnh báo giả (Alert Fatigue) cao; chưa có khả năng hỏi đáp ngôn ngữ tự nhiên và truy vết văn bản gốc. & Đang được sử dụng phổ biến nhưng trải nghiệm người dùng chưa cao, bác sĩ thường tắt bỏ cảnh báo. \\
\bottomrule
\end{tabularx}
\caption{Bảng so sánh và đánh giá các hệ thống CDSS hiện hữu}
\end{table}

Từ bảng so sánh trên, có thể thấy rõ khoảng trống công nghệ hiện nay: các bệnh viện Việt Nam đang thiếu một giải pháp CDSS có chi phí hợp lý, tích hợp sâu tri thức pháp quy nội địa, hoạt động nhẹ nhàng cục bộ và hỗ trợ tương tác thông minh qua ngôn ngữ tự nhiên.
